% !PRINTED 237
\weberpage{237}
%<<CONT
If $a$, $b$, $c$, $\partial$ are four integers without a common divisor
satisfying the condition:
\[
a\partial - bc = n
\]
then one obtains the pairs of numbers $\mu$, $\mu'$ corresponding to one
row by letting, in
\weq{19}{\begin{gathered}
\mu \equiv a\alpha' + c\beta' \\
\mu' \equiv b\alpha' + \partial\beta' \ (\mathrm{mod}\ n)
\end{gathered}}
$\alpha'$, $\beta'$ run through all those values, and only those, for which
$\mu$, $\mu'$ have no common divisor with $n$.

That two pairs of values $\mu$, $\mu'$ corresponding to the congruences (19)
really do belong to the same row follows immediately
from § 65 (3), and it is likewise self-evident that \gesperrt{all pairs of numbers}
of one row are contained in this form, since $\alpha'$, $\beta'$
may be replaced by $h\alpha'$, $h\beta'$ when $h$ is relatively prime to $n$. That
one obtains \gesperrt{all the rows} in this way is shown by
the formulae (10).

\begin{center}
\textbf{§ 69.} \textbf{\gesperrt{The Invariant Equation.}}
\end{center}

Among the transformation equations, one is of particular
interest: the equation whose roots are the $\psi(n)$ quantities
\weq{1}{j\left(\frac{c+\partial\omega}{a}\right)}
or, by the determination (19), § 68, the quantities identical with
them
\weq{2}{j\left(\frac{c+\partial\omega}{a+b\omega}\right)}
where $j(\omega)$ is the invariant defined in § 46.

This equation is called the \gesperrt{invariant equation.}

The function $j(\omega)$ can, by § 54, be represented rationally by $f(\omega)^{24}$,
namely
\weq{3}{j(\omega) = \frac{[f(\omega)^{24} - 16]^3}{f(\omega)^{24}},}
so that, by the results of the preceding section, the
quantities (1) are the roots of an equation whose coefficients
are rational functions of $\varkappa^2$.

But if for $f(\omega)$ in (3) one substitutes one of the developments
found earlier, e.g. § 24, (11):
\[
% [sic] The source prints a stray comma after the lower limit 1.
f(\omega) = q^{-\frac{1}{24}} \prod_{1,}^{\nu} (1+q^{2\nu-1}),
\]
%CONT>>
